\subsection{Model Selection}
Before outlining the methodology for model selection, we motivate the procedure by exploring at the implications of setting certain parameters in an ARMA-Noise model $\bar{\mathbf{w}}_{t\cdot} = \uu + \mathbf{a} $ equal to zero.


\subsubsection{${\text{Case 1: } \uu \text{ is AR(1)}}$}



We have already shown that this case results in $\mathbf{w}$ being an ARMA(1,1) process with innovations $\{b_t\}$ and parameterized by $(\phi,\theta_b,\sigma_b^2)$. Figures \ref{fig:neg_ar_var_effect} and \ref{fig:pos_ar_var_effect} remind us that $\theta_b$ will fall somewhere between $0$ and $-\phi$, converging to the former as $\frac{\sigma_a^2}{\sigma_\eta^2}$ approaches 0 and the latter as it approaches infinity.


Assuming values of $\phi$, $\sigma_a^2$, and $T$ of large enough magnitude to result in significant estimates for $\phi$ and $\theta_b$, and not knowing $\uu$ is AR(1), we use  Equations \ref{eq:H_arman}, \ref{eq:delta}, and \ref{eq:G_implicit} to estimate and perform inference on $\hat{\tau_l}$. In the event $\hat{\theta}$ is insignificant we impose  $\theta = 0$ and re-estimate $\sigma_\eta^2$ using $H_\theta$.
\begin{equation}
    \label{eq:H_theta}
    H_\theta = (1+\theta_b^2)\sigma_b^2 - \sigma_\eta^2 - (1 + \phi^2)\sigma_a^2, \qquad  H_f = \begin{cases}
    \sigma_{\varepsilon*}^2 - \sigma_\varepsilon^2 \\
    \phi_* - \phi \\
    H_\theta
\end{cases}  
\end{equation}

\begin{equation}
    \label{eq:Sigma_o_theta_zero}
    \Sigma_o = \left[\begin{array}{ccc}
        \Sigma_{\phi,\theta_b} & \mathbf{0}  & \mathbf{0} \\
        \mathbf{0}' & 2\sigma_b^4 & 0\\
        \mathbf{0}' & 0 & \frac{2\sigma_\varepsilon^4}{n-1}
    \end{array}\right]
\end{equation}


$$ \tau_o = (\phi,\theta_b,\sigma_b^2, \sigma_\varepsilon^2)', \qquad \tau_l = (\sigma_{\varepsilon*}^2, \phi_*, \sigma_\eta^2)' $$

$$ \frac{\partial H_f}{\partial \tau_o} = \left[\begin{array}{cccc}
    0 & 0 & 0 & -1 \\
    -1 & 0 & 0 & 0 \\
    -2\phi\sigma_a^2 & 2\theta_b\sigma_b^2 & (1 + \theta_b^2) & (1 + \phi^2)n^{-1}
\end{array}\right], \qquad \frac{\partial H_f}{\partial \tau_l} = \left[\begin{array}{ccc}
    1 & 0 & 0  \\
    0 & 1 & 0 \\
    0 & 0 & -1
\end{array}\right] $$

\subsubsection{$\text{Case 2: } \uu \text{ is MA(1)}$}
We now consider the case where $\phi = 0$. If $\uu$ is MA(1), so too is $\bar{\mathbf{w}}_{t\cdot}$, with parameters $\theta_b$ and $\sigma_b^2$. The resulting system of related ACVF's is: 

$$ (1 + \theta_bB)b_t = (1 + \theta B)\eta_t + a_t $$

With this restricted relation we adapt our implicit function:

$$H_\phi = \begin{cases}
    (1 + \theta_b^2)\sigma_b^2 - (1 + \theta^2)\sigma_\eta^2 - \sigma_a^2 \\
    \theta_b\sigma_b^2 - \theta\sigma_\eta^2
\end{cases}$$

Because the added noise in no way alters the lack of auto-regression, fitting an ARMA(1,1) process to $\mathbf{w}$ should yield an insignificant estimate for $\phi$. In this event, $\mathbf{w}$ would be refit with an MA(1) model before estimating $\theta$ and $\sigma_\eta^2$ with an adjusted version of $H$, which we will denote $H_\phi$:

\begin{equation}
    \label{eq:H_phi}
    H_{\phi} = \begin{cases}
        (1 + \theta_b^2)\sigma_b^2 - (1 + \theta^2)\sigma_\eta^2 - \sigma_a^2 \\
        \theta_b\sigma_b^2 - \theta\sigma_\eta^2
    \end{cases} 
\end{equation}

\begin{equation}
    \label{eq:Sigma_o_phi_zero}
    \Sigma_o = \left[ \begin{array}{ccc}
        1-\theta^2 & 0 & 0  \\
        0  & 2\sigma_b^2 & 0 \\
        0 & 0 & \frac{2\sigma_\varepsilon^2}{n-1}
    \end{array} \right]
\end{equation}

$$ \tau_o = (\theta_b, \sigma_b^2, \sigma_\varepsilon^2)' \qquad \tau_l = (\sigma_\varepsilon^2,\theta,\sigma_\eta^2)' $$

$$ H_f = \left[\begin{array}{c}
     \sigma_\varepsilon^2 - \sigma_{\varepsilon*}^2  \\
      H_\phi
\end{array}\right] $$

$$ \frac{\partial H_f}{\partial \tau_o} = \left[ \begin{array}{ccc}
    0 & 0 & -1 \\
    2\theta_b\sigma_b^2 & 1 + \theta_b^2 & -n^{-1} \\
    \sigma_b^2 & \theta_b & 0
\end{array}\right], \qquad \frac{\partial H_f}{\partial \tau_l} = \left[ \begin{array}{ccc}
    1 & 0 & 0 \\
    0 & -2\theta\sigma_\eta^2 & -(1+\theta^2) \\
    0& -\sigma_\eta^2 & -\theta
\end{array}\right] $$


% \subsubsection{$\text{Case 3: } \phi = \theta = 0$}

% It's now worth touching on a relatively uninteresting but nevertheless conceivable case of $\mathbf{u}$ simply being white noise. In this event, $\mathbf{w}$ is also white noise, with variance $\sigma_b^2 = \sigma_\eta^2 + \sigma_a^2$. If fitting an ARMA(1,1) model to $\mathbf{w}$ yielded insignificant estimates for both $\phi$ and $\theta$, we would first refit with an MA(1) and then a white noise model if the resulting MA estimate was insignificant.

% We then estimate $\sigma_\eta^2$ with $\hat{\sigma}_b^2 - \hat{\sigma}_a^2$ and generate confidence intervals based on:

% $$\sqrt{T}(\hat{\sigma}_\eta^2 - \sigma_\eta^2) \xrightarrow{d}N\left[0,2\sigma_b^4 - \frac{2\sigma_\varepsilon^4}{n^2(n-1)}\right]$$

% \textcolor{blue}{check the above expression}
\subsubsection{$\text{Case 3: } \theta_b = 0$}
We now turn to a case that was hinted at by Figures \ref{fig:neg_ar_var_effect} and \ref{fig:pos_ar_var_effect}, namely that $\bar{\mathbf{w}}$ is AR(1). We've already established  the sum of MA(1) and white noise is still MA(1). This means that $\uu$ must be ARMA(1,1), setting up the following system of ACVF's:

$$ (1-\phi B)^{-1}b_t = (1-\phi B)^{-1}(1 + \theta B)\eta_t + a_t $$

$$ b_t = (1 + \theta B)\eta_t + (1 -\phi B)a_t $$

We calculate respective ACVF's at lags 0 and 1 to obtain:

\begin{equation}
    \label{eq:H_thetab}
    H_{\theta_b} =  \begin{cases}
    \sigma_b^2 - (1 + \theta^2)\sigma_\eta^2 - (1 + \phi^2)\sigma_a^2 \\
    \theta \sigma_\eta^2 - \phi \sigma_a^2
    \end{cases}
\end{equation}


Note that the second equation provides a critical piece of information, namely that $\theta_b$ being zero necessitates that $\frac{\theta}{\phi} = \frac{\sigma^2_{a}}{\sigma^2_{\eta}}$. This gives rise to the following:


\begin{equation}
    \label{eq:Sigma_o_thetab}
    \Sigma_o = \left[\begin{array}{ccc}
        1 - \phi^2 & 0 & 0 \\
        0 & 2\sigma_b^4 & 0 \\
        0 & 0 & \frac{2\sigma_\varepsilon^4}{n-1}
    \end{array}\right]
\end{equation}

$$ \tau_o = (\phi,\sigma_b^2, \sigma_\varepsilon^2)', \qquad \tau_l = (\sigma_{\varepsilon*}^2, \phi_*, \theta, \sigma_\eta^2), \qquad H_f = \begin{cases}
    \sigma_{\varepsilon*}^2 - \sigma_{\varepsilon}^2 \\
    \phi_* - \phi \\
    H_{\theta_b}
\end{cases} $$



$$ \frac{\partial H_f}{\partial \tau_o} = \left[ \begin{array}{ccc}
    0 & 0 & -1 \\
    -1 & 0 & 0 \\
    -2\phi\sigma_a^2 & 1 & -(1 + \phi^2)\hat{m} \\
    -\sigma_a^2 & 0 & -\phi\hat{m}
\end{array}\right], \qquad \frac{\partial H_f}{\partial \tau_l} = \left[ \begin{array}{cccc}
    1 & 0 & 0 & 0 \\
    0 & 1 & 0 & 0 \\
    0 & 0 & -2\theta\sigma_\eta^2 & -(1+\theta^2) \\
    0& 0 &  \sigma_\eta^2 & \theta
\end{array}\right] $$


In the event of an insignificant estimate for $\theta_b$, we would refit $\mathbf{w}$ with an AR(1) process and then use $H_{\theta_b}$ to fit the remaining parameters. We synthesize the implications of these cases with the following procedure for model selection, which is performed using the parameter estimates of the converged REMARCS model.


\begin{figure}[htbp]
    \centering
    \vspace{.5in}
\resizebox{\textwidth}{!}{
\begin{tikzpicture}[scale=1.2, every node/.style={transform shape}, node distance=2.5cm]{}
\small
    % Nodes vertically aligned
    \node (start) [startstop, align=center] {Start \\
    $\yy = X_t\uu + \bm{\varepsilon}$ \\
    $\mathbf{w} = \uu + \mathbf{a}$  \\
    $\mathbf{x} = \yy - X_t\mathbf{w}$};
    % \node (model) [process, right = 2.5cm of start, align=center]
    % {};
    \node (noisehat) [process, right=3.5cm of start]{$1)\widehat{\sigma_\varepsilon^2} = \frac{||\mathbf{x}||^2}{N-T}$};
    \node (observed) [process, below=2.5cm of noisehat,align=center] 
        {Fit $\mathbf{w}$ with ARMA(1,1)\\
        to get $\hat{\phi}, \widehat{\theta_b}, \widehat{\sigma_b^2}$};
    \node (latent) [process, below = 2.75cm of observed] 
        {$2)\hat{\theta}, \widehat{\sigma_\eta^2} \text{ are roots of }$\ref{eq:H_arman}};
    \node (phizero) [process, left = 2.5cm of observed, align=center] 
         {$\text{Set }\hat{\phi} = 0$ and\\
          refit $\mathbf{w}$ with MA(1)\\
          to update $\hat{\theta_b}$ and $\widehat{\sigma_b^2}$
         };
    \node (thetabzero) [process, right = 2.5cm of observed,align=center] 
         {$4)\text{Set }\widehat{\theta_b} = 0$ and \\
         refit $\mathbf{w}$ with an AR(1) \\
         to update $\hat{\phi}$ and $\widehat{\sigma_b^2}$};
    \node (thetabzerotheta) [process, below = 2.5cm of thetabzero] 
         {$5)\widehat{\theta}, \hat{\sigma}_\eta^2 \text{ are roots of }$\ref{eq:H_thetab}};
    \node (phizerotheta) [process, below = 2.5cm of phizero] 
         {$5)\widehat{\theta}, \hat{\sigma}_\eta^2 \text{ are roots of }$\ref{eq:H_phi}};

    \node (done) [startstop, below = 2.5cm of thetabzerotheta] {Done};


    % Arrows
    \draw [arrow] (start) -- (noisehat)
        node[pos=0.3, anchor = south]{};

    \draw [arrow] (noisehat) -- (observed);
    \draw [arrow] (observed) -- (phizero) 
        node[pos=.5,anchor = south]{If $\hat{\phi}$ is}
        node[pos=.5,anchor=north]{insignificant};
        \draw [arrow] (observed) -- (thetabzero) 
        node[pos=.5,anchor = south]{If only $\hat{\theta}_b$ }
        node[pos=.5,anchor=north]{is insignificant};
    \draw [arrow] (thetabzero) -- (thetabzerotheta);
        \draw [arrow] (phizero) -- (phizerotheta);

    \draw [arrow] (observed) -- (latent)
    node[pos=.5, anchor=west]{else};
    \draw [arrow] (thetabzerotheta) -- (done)
    node[pos=.5, anchor=west]{};
    \draw [arrow] (latent) |- (done)
    node[pos=.5, anchor=west]{};
    \draw [arrow] (phizerotheta) |- (done)
    node[pos=.5, anchor=west]{};


\end{tikzpicture}
}
    \caption{Decision tree for selecting parameters to include in ARMA-Noise model}
    \label{fig:modelselection}
\end{figure}

